% Options for packages loaded elsewhere
% Options for packages loaded elsewhere
\PassOptionsToPackage{unicode}{hyperref}
\PassOptionsToPackage{hyphens}{url}
\PassOptionsToPackage{dvipsnames,svgnames,x11names}{xcolor}
%
\documentclass[
  russian,
]{article}
\usepackage{xcolor}
\usepackage[top=20mm,left=20mm,right=20mm,bottom=20mm]{geometry}
\usepackage{amsmath,amssymb}
\setcounter{secnumdepth}{5}
\usepackage{iftex}
\ifPDFTeX
  \usepackage[T1]{fontenc}
  \usepackage[utf8]{inputenc}
  \usepackage{textcomp} % provide euro and other symbols
\else % if luatex or xetex
  \usepackage{unicode-math} % this also loads fontspec
  \defaultfontfeatures{Scale=MatchLowercase}
  \defaultfontfeatures[\rmfamily]{Ligatures=TeX,Scale=1}
\fi
\usepackage{lmodern}
\ifPDFTeX\else
  % xetex/luatex font selection
  \setmainfont[]{Times New Roman}
  \setsansfont[]{Arial}
  \setmonofont[]{Courier New}
\fi
% Use upquote if available, for straight quotes in verbatim environments
\IfFileExists{upquote.sty}{\usepackage{upquote}}{}
\IfFileExists{microtype.sty}{% use microtype if available
  \usepackage[]{microtype}
  \UseMicrotypeSet[protrusion]{basicmath} % disable protrusion for tt fonts
}{}
\makeatletter
\@ifundefined{KOMAClassName}{% if non-KOMA class
  \IfFileExists{parskip.sty}{%
    \usepackage{parskip}
  }{% else
    \setlength{\parindent}{0pt}
    \setlength{\parskip}{6pt plus 2pt minus 1pt}}
}{% if KOMA class
  \KOMAoptions{parskip=half}}
\makeatother
% Make \paragraph and \subparagraph free-standing
\makeatletter
\ifx\paragraph\undefined\else
  \let\oldparagraph\paragraph
  \renewcommand{\paragraph}{
    \@ifstar
      \xxxParagraphStar
      \xxxParagraphNoStar
  }
  \newcommand{\xxxParagraphStar}[1]{\oldparagraph*{#1}\mbox{}}
  \newcommand{\xxxParagraphNoStar}[1]{\oldparagraph{#1}\mbox{}}
\fi
\ifx\subparagraph\undefined\else
  \let\oldsubparagraph\subparagraph
  \renewcommand{\subparagraph}{
    \@ifstar
      \xxxSubParagraphStar
      \xxxSubParagraphNoStar
  }
  \newcommand{\xxxSubParagraphStar}[1]{\oldsubparagraph*{#1}\mbox{}}
  \newcommand{\xxxSubParagraphNoStar}[1]{\oldsubparagraph{#1}\mbox{}}
\fi
\makeatother


\usepackage{longtable,booktabs,array}
\newcounter{none} % for unnumbered tables
\usepackage{calc} % for calculating minipage widths
% Correct order of tables after \paragraph or \subparagraph
\usepackage{etoolbox}
\makeatletter
\patchcmd\longtable{\par}{\if@noskipsec\mbox{}\fi\par}{}{}
\makeatother
% Allow footnotes in longtable head/foot
\IfFileExists{footnotehyper.sty}{\usepackage{footnotehyper}}{\usepackage{footnote}}
\makesavenoteenv{longtable}
\usepackage{graphicx}
\makeatletter
\newsavebox\pandoc@box
\newcommand*\pandocbounded[1]{% scales image to fit in text height/width
  \sbox\pandoc@box{#1}%
  \Gscale@div\@tempa{\textheight}{\dimexpr\ht\pandoc@box+\dp\pandoc@box\relax}%
  \Gscale@div\@tempb{\linewidth}{\wd\pandoc@box}%
  \ifdim\@tempb\p@<\@tempa\p@\let\@tempa\@tempb\fi% select the smaller of both
  \ifdim\@tempa\p@<\p@\scalebox{\@tempa}{\usebox\pandoc@box}%
  \else\usebox{\pandoc@box}%
  \fi%
}
% Set default figure placement to htbp
\def\fps@figure{htbp}
\makeatother



\ifLuaTeX
\usepackage[bidi=basic,shorthands=off,provide=*]{babel}
\else
\usepackage[bidi=default,shorthands=off,provide=*]{babel}
\fi
\ifPDFTeX
\else
\babelfont{rm}[]{Times New Roman}
\fi
\ifLuaTeX
  \usepackage{selnolig} % disable illegal ligatures
\fi


\setlength{\emergencystretch}{3em} % prevent overfull lines

\providecommand{\tightlist}{%
  \setlength{\itemsep}{0pt}\setlength{\parskip}{0pt}}



 


\makeatletter
\@ifpackageloaded{tcolorbox}{}{\usepackage[skins,breakable]{tcolorbox}}
\@ifpackageloaded{fontawesome5}{}{\usepackage{fontawesome5}}
\definecolor{quarto-callout-color}{HTML}{909090}
\definecolor{quarto-callout-note-color}{HTML}{0758E5}
\definecolor{quarto-callout-important-color}{HTML}{CC1914}
\definecolor{quarto-callout-warning-color}{HTML}{EB9113}
\definecolor{quarto-callout-tip-color}{HTML}{00A047}
\definecolor{quarto-callout-caution-color}{HTML}{FC5300}
\definecolor{quarto-callout-color-frame}{HTML}{acacac}
\definecolor{quarto-callout-note-color-frame}{HTML}{4582ec}
\definecolor{quarto-callout-important-color-frame}{HTML}{d9534f}
\definecolor{quarto-callout-warning-color-frame}{HTML}{f0ad4e}
\definecolor{quarto-callout-tip-color-frame}{HTML}{02b875}
\definecolor{quarto-callout-caution-color-frame}{HTML}{fd7e14}
\makeatother
\makeatletter
\@ifpackageloaded{caption}{}{\usepackage{caption}}
\AtBeginDocument{%
\ifdefined\contentsname
  \renewcommand*\contentsname{Содержание}
\else
  \newcommand\contentsname{Содержание}
\fi
\ifdefined\listfigurename
  \renewcommand*\listfigurename{Список Иллюстраций}
\else
  \newcommand\listfigurename{Список Иллюстраций}
\fi
\ifdefined\listtablename
  \renewcommand*\listtablename{Список Таблиц}
\else
  \newcommand\listtablename{Список Таблиц}
\fi
\ifdefined\figurename
  \renewcommand*\figurename{Рисунок}
\else
  \newcommand\figurename{Рисунок}
\fi
\ifdefined\tablename
  \renewcommand*\tablename{Таблица}
\else
  \newcommand\tablename{Таблица}
\fi
}
\@ifpackageloaded{float}{}{\usepackage{float}}
\floatstyle{ruled}
\@ifundefined{c@chapter}{\newfloat{codelisting}{h}{lop}}{\newfloat{codelisting}{h}{lop}[chapter]}
\floatname{codelisting}{Список}
\newcommand*\listoflistings{\listof{codelisting}{Список Каталогов}}
\makeatother
\makeatletter
\makeatother
\makeatletter
\@ifpackageloaded{caption}{}{\usepackage{caption}}
\@ifpackageloaded{subcaption}{}{\usepackage{subcaption}}
\makeatother
\usepackage{bookmark}
\IfFileExists{xurl.sty}{\usepackage{xurl}}{} % add URL line breaks if available
\urlstyle{same}
% fallback for those not using the hyperref driver hyperxmp:
\makeatletter
\@ifundefined{xmpquote}{\newcommand{\xmpquote}[1]{#1}}{}
\makeatother
\hypersetup{
  pdftitle={Задачи по Эконометрике-1: Прогнозирование},
  pdfauthor={Н. В. Артамонов, А. Н. Лата (МГИМО МИД России)},
  pdflang={ru},
  colorlinks=true,
  linkcolor={blue},
  filecolor={Maroon},
  citecolor={Blue},
  urlcolor={Blue},
  pdfcreator={LaTeX via pandoc}}


\title{Задачи по Эконометрике-1: Прогнозирование}
\author{Н. В. Артамонов, А. Н. Лата (МГИМО МИД России)}
\date{}
\begin{document}
\maketitle

\renewcommand*\contentsname{Содержание}
{
\setcounter{tocdepth}{3}
\tableofcontents
}

\begin{tcolorbox}[enhanced jigsaw, arc=.35mm, bottomrule=.15mm, bottomtitle=1mm, breakable, colback=white, colbacktitle=quarto-callout-important-color!10!white, colframe=quarto-callout-important-color-frame, coltitle=black, left=2mm, leftrule=.75mm, opacityback=0, opacitybacktitle=0.6, rightrule=.15mm, title=\textcolor{quarto-callout-important-color}{\faExclamation}\hspace{0.5em}{Важно}, titlerule=0mm, toprule=.15mm, toptitle=1mm]

\begin{itemize}
\tightlist
\item
  Логарифм натуральный; константа включается в модель по умолчанию.
\item
  Данные лежат в папке \texttt{datasets}, описание переменных --- в
  \texttt{datasets/data-description.md}.
\item
  Ответы округляются до 2-х десятичных знаков.
\item
  Коэффициенты подписаны так, как их печатает statsmodels:
  \texttt{Intercept} --- константа, \texttt{np.log(x)} --- \(\log x\),
  \texttt{I(x\ **\ 2)} --- \(x^2\).
\item
  Прогноз считается по неокруглённым оценкам (метод \texttt{predict});
  расчёт по округлённым числам из таблицы может дать другой ответ.
\item
  В логарифмических моделях прогноз в уровнях --- экспонента от прогноза
  логарифма (без поправки).
\end{itemize}

\end{tcolorbox}

\section{Прогноз в модели в
уровнях}\label{ux43fux440ux43eux433ux43dux43eux437-ux432-ux43cux43eux434ux435ux43bux438-ux432-ux443ux440ux43eux432ux43dux44fux445}

\subsection{sleep equation \#1}\label{sleep-equation-1}

Для набора данных \texttt{sleep75} (\texttt{sleep} --- время ночного
сна, \texttt{totwrk} --- время работы, обе в минутах в неделю)
рассмотрим регрессию

\texttt{sleep\ \textasciitilde{}\ totwrk\ +\ age\ +\ south\ +\ male\ +\ smsa\ +\ yngkid\ +\ marr}

Результаты оценивания:

{\def\LTcaptype{none} % do not increment counter
\begin{longtable}[]{@{}lr@{}}
\toprule\noalign{}
Переменная & Оценка \\
\midrule\noalign{}
\endhead
\bottomrule\noalign{}
\endlastfoot
Intercept & 3451 \\
totwrk & -0.1692 \\
age & 2.689 \\
south & 101.6 \\
male & 87.67 \\
smsa & -54.75 \\
yngkid & -13.96 \\
marr & 31.21 \\
\end{longtable}
}

Модель была подогнана по 706 наблюдениям.

Рассмотрим трёх людей с характеристиками

{\def\LTcaptype{none} % do not increment counter
\begin{longtable}[]{@{}rrrrrrrr@{}}
\toprule\noalign{}
Человек & totwrk & age & south & male & smsa & yngkid & marr \\
\midrule\noalign{}
\endhead
\bottomrule\noalign{}
\endlastfoot
1 & 2150 & 37 & 0 & 1 & 1 & 0 & 1 \\
2 & 1950 & 28 & 1 & 1 & 0 & 1 & 0 \\
3 & 2240 & 26 & 0 & 0 & 1 & 0 & 0 \\
\end{longtable}
}

Вычислите прогноз для каждого человека по подогнанной модели.
\textbf{Ответ округлите до 2-х десятичных знаков.}

\textbf{Прогнозы sleep (минут в неделю):}

{\def\LTcaptype{none} % do not increment counter
\begin{longtable}[]{@{}rr@{}}
\toprule\noalign{}
Человек & Прогноз \\
\midrule\noalign{}
\endhead
\bottomrule\noalign{}
\endlastfoot
1 & 3250.68 \\
2 & 3371.46 \\
3 & 3086.98 \\
\end{longtable}
}

\subsection{sleep equation \#2}\label{sleep-equation-2}

Для набора данных \texttt{sleep75} рассмотрим регрессию

\texttt{sleep\ \textasciitilde{}\ totwrk\ +\ age\ +\ I(age**2)\ +\ south\ +\ male}

Результаты оценивания:

{\def\LTcaptype{none} % do not increment counter
\begin{longtable}[]{@{}lr@{}}
\toprule\noalign{}
Переменная & Оценка \\
\midrule\noalign{}
\endhead
\bottomrule\noalign{}
\endlastfoot
Intercept & 3614 \\
totwrk & -0.1673 \\
age & -6.255 \\
I(age ** 2) & 0.1089 \\
south & 114.5 \\
male & 90.46 \\
\end{longtable}
}

Модель была подогнана по 706 наблюдениям.

Рассмотрим трёх людей с характеристиками

{\def\LTcaptype{none} % do not increment counter
\begin{longtable}[]{@{}rrrrr@{}}
\toprule\noalign{}
Человек & totwrk & age & south & male \\
\midrule\noalign{}
\endhead
\bottomrule\noalign{}
\endlastfoot
1 & 2160 & 32 & 1 & 0 \\
2 & 1720 & 24 & 0 & 1 \\
3 & 2390 & 44 & 0 & 1 \\
\end{longtable}
}

Вычислите прогноз для каждого человека по подогнанной модели.
\textbf{Ответ округлите до 2-х десятичных знаков.}

\textbf{Прогнозы sleep (минут в неделю):}

{\def\LTcaptype{none} % do not increment counter
\begin{longtable}[]{@{}rr@{}}
\toprule\noalign{}
Человек & Прогноз \\
\midrule\noalign{}
\endhead
\bottomrule\noalign{}
\endlastfoot
1 & 3278.46 \\
2 & 3329.25 \\
3 & 3240.11 \\
\end{longtable}
}

\section{Прогноз в логарифмической
модели}\label{ux43fux440ux43eux433ux43dux43eux437-ux432-ux43bux43eux433ux430ux440ux438ux444ux43cux438ux447ux435ux441ux43aux43eux439-ux43cux43eux434ux435ux43bux438}

\subsection{wage equation \#1}\label{wage-equation-1}

Для набора данных \texttt{wage2} (\texttt{wage} --- месячный заработок)
рассмотрим регрессию

\texttt{np.log(wage)\ \textasciitilde{}\ age\ +\ IQ\ +\ south\ +\ married\ +\ urban}

Результаты оценивания:

{\def\LTcaptype{none} % do not increment counter
\begin{longtable}[]{@{}lr@{}}
\toprule\noalign{}
Переменная & Оценка \\
\midrule\noalign{}
\endhead
\bottomrule\noalign{}
\endlastfoot
Intercept & 4.974 \\
age & 0.02125 \\
IQ & 0.008201 \\
south & -0.09898 \\
married & 0.201 \\
urban & 0.175 \\
\end{longtable}
}

Модель была подогнана по 935 наблюдениям.

Рассмотрим трёх людей с характеристиками

{\def\LTcaptype{none} % do not increment counter
\begin{longtable}[]{@{}rrrrrr@{}}
\toprule\noalign{}
Человек & age & IQ & south & married & urban \\
\midrule\noalign{}
\endhead
\bottomrule\noalign{}
\endlastfoot
1 & 36 & 105 & 1 & 1 & 1 \\
2 & 29 & 123 & 0 & 1 & 0 \\
3 & 25 & 112 & 1 & 0 & 1 \\
\end{longtable}
}

Вычислите прогноз \textbf{зарплаты} для каждого человека по подогнанной
модели. Прогноз в уровнях получите как экспоненту от прогноза логарифма
(без поправки). \textbf{Ответ округлите до 2-х десятичных знаков.}

\textbf{Прогнозы wage (месячный заработок):}

{\def\LTcaptype{none} % do not increment counter
\begin{longtable}[]{@{}rrr@{}}
\toprule\noalign{}
Человек & Прогноз np.log(wage) & Прогноз wage \\
\midrule\noalign{}
\endhead
\bottomrule\noalign{}
\endlastfoot
1 & 6.877 & 969.95 \\
2 & 6.800 & 897.93 \\
3 & 6.500 & 665.07 \\
\end{longtable}
}

\subsection{wage equation \#2}\label{wage-equation-2}

Для набора данных \texttt{wage1} (\texttt{wage} --- средний почасовой
заработок) рассмотрим регрессию

\texttt{np.log(wage)\ \textasciitilde{}\ exper\ +\ I(exper**2)\ +\ female\ +\ married\ +\ smsa}

Результаты оценивания:

{\def\LTcaptype{none} % do not increment counter
\begin{longtable}[]{@{}lr@{}}
\toprule\noalign{}
Переменная & Оценка \\
\midrule\noalign{}
\endhead
\bottomrule\noalign{}
\endlastfoot
Intercept & 1.25 \\
exper & 0.03688 \\
I(exper ** 2) & -0.0007717 \\
female & -0.3627 \\
married & 0.1438 \\
smsa & 0.2726 \\
\end{longtable}
}

Модель была подогнана по 526 наблюдениям.

Рассмотрим трёх людей с характеристиками

{\def\LTcaptype{none} % do not increment counter
\begin{longtable}[]{@{}rrrrr@{}}
\toprule\noalign{}
Человек & exper & female & married & smsa \\
\midrule\noalign{}
\endhead
\bottomrule\noalign{}
\endlastfoot
1 & 5 & 1 & 1 & 1 \\
2 & 26 & 0 & 0 & 1 \\
3 & 38 & 1 & 1 & 0 \\
\end{longtable}
}

Вычислите прогноз \textbf{зарплаты} для каждого человека по подогнанной
модели. Прогноз в уровнях получите как экспоненту от прогноза логарифма
(без поправки). \textbf{Ответ округлите до 2-х десятичных знаков.}

\textbf{Прогнозы wage (средний почасовой заработок):}

{\def\LTcaptype{none} % do not increment counter
\begin{longtable}[]{@{}rrr@{}}
\toprule\noalign{}
Человек & Прогноз np.log(wage) & Прогноз wage \\
\midrule\noalign{}
\endhead
\bottomrule\noalign{}
\endlastfoot
1 & 1.469 & 4.35 \\
2 & 1.960 & 7.10 \\
3 & 1.319 & 3.74 \\
\end{longtable}
}

\subsection{output equation \#1}\label{output-equation-1}

Для набора данных \texttt{Labour} (\texttt{output} --- добавленная
стоимость, \texttt{capital} --- основной капитал, обе в млн евро;
\texttt{labour} --- число работников) рассмотрим регрессию
\textbf{np.log(output) на np.log(capital), np.log(labour)}.

Результаты оценивания:

{\def\LTcaptype{none} % do not increment counter
\begin{longtable}[]{@{}lr@{}}
\toprule\noalign{}
Переменная & Оценка \\
\midrule\noalign{}
\endhead
\bottomrule\noalign{}
\endlastfoot
Intercept & -1.711 \\
np.log(capital) & 0.2076 \\
np.log(labour) & 0.7148 \\
\end{longtable}
}

Модель была подогнана по 569 наблюдениям.

Рассмотрим три фирмы с характеристиками

{\def\LTcaptype{none} % do not increment counter
\begin{longtable}[]{@{}rrr@{}}
\toprule\noalign{}
Фирма & capital & labour \\
\midrule\noalign{}
\endhead
\bottomrule\noalign{}
\endlastfoot
1 & 2.97 & 85 \\
2 & 10.45 & 60 \\
3 & 3.85 & 105 \\
\end{longtable}
}

Вычислите прогноз \textbf{выпуска} для каждой фирмы по подогнанной
модели. Прогноз в уровнях получите как экспоненту от прогноза логарифма
(без поправки). \textbf{Ответ округлите до 2-х десятичных знаков.}

\textbf{Прогнозы output (млн евро):}

{\def\LTcaptype{none} % do not increment counter
\begin{longtable}[]{@{}rrr@{}}
\toprule\noalign{}
Фирма & Прогноз np.log(output) & Прогноз output \\
\midrule\noalign{}
\endhead
\bottomrule\noalign{}
\endlastfoot
1 & 1.690 & 5.42 \\
2 & 1.702 & 5.49 \\
3 & 1.895 & 6.65 \\
\end{longtable}
}

\subsection{output equation \#2}\label{output-equation-2}

Для набора данных \texttt{Labour} рассмотрим регрессию

\begin{verbatim}
np.log(output) ~ np.log(capital) + np.log(labour) +
    I(np.log(capital)**2) + I(np.log(labour)**2)
\end{verbatim}

Результаты оценивания:

{\def\LTcaptype{none} % do not increment counter
\begin{longtable}[]{@{}lr@{}}
\toprule\noalign{}
Переменная & Оценка \\
\midrule\noalign{}
\endhead
\bottomrule\noalign{}
\endlastfoot
Intercept & -1.304 \\
np.log(capital) & 0.1831 \\
np.log(labour) & 0.5153 \\
I(np.log(capital) ** 2) & 0.02275 \\
I(np.log(labour) ** 2) & 0.02026 \\
\end{longtable}
}

Модель была подогнана по 569 наблюдениям.

Рассмотрим три фирмы с характеристиками

{\def\LTcaptype{none} % do not increment counter
\begin{longtable}[]{@{}rrr@{}}
\toprule\noalign{}
Фирма & capital & labour \\
\midrule\noalign{}
\endhead
\bottomrule\noalign{}
\endlastfoot
1 & 22.14 & 407 \\
2 & 7.32 & 197 \\
3 & 0.67 & 31 \\
\end{longtable}
}

Вычислите прогноз \textbf{выпуска} для каждой фирмы по подогнанной
модели. Прогноз в уровнях получите как экспоненту от прогноза логарифма
(без поправки). \textbf{Ответ округлите до 2-х десятичных знаков.}

\textbf{Прогнозы output (млн евро):}

{\def\LTcaptype{none} % do not increment counter
\begin{longtable}[]{@{}rrr@{}}
\toprule\noalign{}
Фирма & Прогноз np.log(output) & Прогноз output \\
\midrule\noalign{}
\endhead
\bottomrule\noalign{}
\endlastfoot
1 & 3.309 & 27.37 \\
2 & 2.439 & 11.46 \\
3 & 0.635 & 1.89 \\
\end{longtable}
}




\end{document}
